\documentclass[a4paper]{article}
\usepackage{amsmath, amssymb}
\usepackage[swedish]{babel}
\usepackage{graphicx}
\graphicspath{{imgs/}}
\usepackage{float}
\usepackage[hidelinks]{hyperref}
\usepackage{fancyhdr}
\usepackage[parfill]{parskip}
\usepackage[margin=1in]{geometry}
\usepackage{listings}
\usepackage{xcolor}
\setcounter{secnumdepth}{0}
\pagestyle{fancy}
\fancyhf{}
\setlength{\headheight}{12.31253pt}
\newcommand{\getauthor}{Eric Johansson (erjo2002@student.miun.se)} %Author
\newcommand{\gettitle}{Laboration 1 \\ Ekvationer och differensformler} %Title
\newcommand{\getcourse}{(MA069G, Matematisk Modellering, 6hp)} %Course
\newcommand{\getsupervisor}{Cornelia Schiebold \\Magnus Eriksson}


\definecolor{codegreen}{rgb}{0,0.6,0}
\definecolor{codegray}{rgb}{0.5,0.5,0.5}
\definecolor{codepurple}{rgb}{0.58,0,0.82}
\definecolor{backcolour}{rgb}{0.95,0.95,0.92}

\lstdefinestyle{mystyle}{
    backgroundcolor=\color{backcolour},   
    commentstyle=\color{codegreen},
    keywordstyle=\color{magenta},
    numberstyle=\tiny\color{codegray},
    stringstyle=\color{codepurple},
    basicstyle=\ttfamily\footnotesize,
    breakatwhitespace=false,         
    breaklines=true,                 
    captionpos=b,                    
    keepspaces=true,                 
    numbers=left,                    
    numbersep=5pt,                  
    showspaces=false,                
    showstringspaces=false,
    showtabs=false,                  
    tabsize=2
}

\lstset{style=mystyle}

\begin{document}

\begin{titlepage}
    \begin{center}
        \vspace*{1cm}
        \includegraphics[width=0.8\textwidth]{msu.png}\\[0.5cm]
        \Large
        Institution of Information Systems and -Technology (IST)\\[1cm]
        \Huge
        \textbf{\gettitle}

        \large
        \getcourse{}

        \vspace{1cm}

        \Large
        \textbf{\getauthor}\\
        \vfill
        \vspace{0.8cm}
        \small
        \today \\
        \Large
        \textbf{Handledare:}\\
        \getsupervisor{}
    \end{center}
\end{titlepage}

\lhead{Laboration 1}
\tableofcontents
\newpage
    
Denna laboration ger er möjlighet att bekanta er med hur man använder 
Python för att lösa några enkla ickelinjära ekvationer samt undersöka 
några finita differensformler.

För att utföra laborationen rekommenderas det att biblioteken numpy,
sympy och matplotlib används. De installeras genom att skriva ``pip install \textit{namn}'' i terminalen.
För många vanliga matematiska funktioner som till exempel sinus, kvadratroten och log krävs att ni importerar
math-biblioteket i er skript.

Ett exempel på start av er fil kan se ut såhär:
\begin{lstlisting}[language=Python]
"""
Laboration X, uppgift Y 
namn1, namn2, namn3
Datum
"""
import math 
import numpy as np
import sympy as sp 
from matplotlib import pyplot
\end{lstlisting}

Därefter kan ni använda funktionaliteten från de olika biblioteken genom att 
kalla deras funktioner. Exempelvis math.sqrt(), math.cos(), sp.symbos()

I rapporten skall ni bifoga er kod, beräkningsresultat, diagram, korta svar och förklaringar. 
Allas namn och datum måste ingå i dokumentet.

\section{Uppgift 1}

Betrakta funktionen \( f(x) \) enligt :
\begin{equation} \label{eqn:fxq1}
f(x)=\frac{\sqrt{(x-5)^5}+2\cos (\pi\sqrt{x})}{\sqrt{x+4 \ln (x-\pi)}-1}
\end{equation}
Om vi vill beräkna \( f'(x) \) i en viss punkt \( x=x_0 \) 
(i detta fall är \( x_0=7 \)) kan vi använda differensformlerna

\begin{equation} \label{eqn:framdiff}
f'(x)=\frac{f(x+h)-f(x)}{h}+E_1(h)=D_1(x,h)+E_1(h)
\end{equation}
\begin{equation} \label{eqn:centerdiff}
f'(x)=\frac{f(x+h)-f(x-h)}{2h}+E_2(h)=D_2(x,h)+E_2(h)
\end{equation}

där (\ref{eqn:framdiff}) är framåtdifferens och (\ref{eqn:centerdiff}) är centraldifferens. 

Skriv ett skript som berkänar och skriver ut \( D_1(x_0,h) \)
och  \( D_2(x_0,h) \) för \( h=0.04,0.02,0,01 \). 

Ett mycket noggrant resultat som är korrekt till 15 decimaler är 

\begin{equation} \label{eqn:fxexakt}
f'(x_0)\approx 2.164821222303575
\end{equation}

\textbf{Svara på följande frågor i rapporten}\\
Vilken av formlerna  \( D_1(x_0,h) \) och  \( D_2(x_0,h) \) var noggrannast, 
d.v.s skilde sig minst från värdet i (\ref{eqn:fxexakt})? Varför är den formeln mer noggrann?


\newpage
\section{Uppgift 2} 
Om man skulle testa att manuellt bestämma \( f'(x) \) av (\ref{eqn:fxq1})
skulle man inse att det blir rätt så bökigt. Däremot kan man använda sig av Sympy för att räkna det symboliskt.

Skriv först:
\begin{lstlisting}[language=Python]
import sympy as sy

x = sy.symbols("x")
f = (sy.sqrt((x-5)**5) + 2*sy.cos(sy.pi*sy.sqrt(x)))/(sy.sqrt(x+4*sy.log(x-sy.pi))-1))
\end{lstlisting}
vilket anger \( f(x) \) symboliskt i python. Därefter skriv:

\begin{lstlisting}[language=Python]
  f_prim = f.diff()
\end{lstlisting}
vilket ger uttrycket för derivatan. 

Sedan, skriv:

\begin{lstlisting}[language=Python]
  f_prim_x0=f_prim.subs(x,7)
\end{lstlisting}
Vilket byter ut x till 7. För att sedan se resultatet skriv:

\begin{lstlisting}[language=Python]
  print(f_prim_x0.evalf(16))
\end{lstlisting}
Detta kommer skriva ut svaret med 16 decimaler och ni kan se 
att det (i princip) exakta värdet på derivatan är 2.164821222303575


\textbf{Svara på följande frågor i rapporten}\\
Vad betyder det att räkna symboliskt? 
Vad kallas det när vi räknar som vanligt som i Uppgift 1?

\newpage
\section{Uppgift 3}

Med hjälp av Taylorutveckling kan man visa att framåtdifferensen i (\ref{eqn:framdiff}) har 
en trunkeringsserie \( E_1(h)\approx k_1h(O(h)) \)(jämfört med centralformeln i (\ref{eqn:centerdiff}) 
som har \(E_1(h)\approx k_1h_2(O(h^2))\), där konstanterna \(k_1\)  och \(c_1\)  är oberoende av \(h\).

Om vi då tar framåtdifferensen i (\ref{eqn:framdiff}) får vi:
\begin{equation} \label{eqn:framdiff_k}
f'(x_0)\approx \frac{f(x_0+h)-f(x_0)}{h}+k_1h
\end{equation}

Om vi istället använder halva stegkängden, \( \frac{h}{2} \) , får vi:
\begin{equation} \label{eqn:framdiff_halva_h}
f'(x_0)\approx \frac{f(x_0+\frac{h}{2})-f(x_0)}{\frac{h}{2}}+k_1 \frac{h}{2}
\end{equation}

Differensen mellan (\ref{eqn:framdiff_k}) och (\ref{eqn:framdiff_halva_h}) ger att vi kan bestämma \( k_1 \) till:

\begin{equation} \label{eqn:k}
k_1=\frac{2\left(2f\left(x_0+\frac{h}{2}\right)-f\left(x_0+h\right)-f\left(x_0\right)\right)}{h^2}
\end{equation}

Använd \( h=0.02 \) i (\ref{eqn:k}) och bestäm ett värde på \( k_1 \). Använd sedan värdet på \( k_1 \) i (\ref{eqn:framdiff_k})
för att bestämma ett värde på \( f'(x_0) \)

\textbf{Svara på följande frågor i rapporten}
Varför blir inte resultatet exakt fastän vi beräknat ut \( k_1 \)?

\section{Uppgift 4}
Det finns även differensformler för andraderivatan: (Centraldifferensen för andraderivatan)
\begin{equation} \label{eqn:fbisscenter}
    f''(x)=\frac{f(x-h)-2f(x)+f(x+h)}{h^2}+E_3(h)=D_3(x,h)+E_3(h)
\end{equation}

Beräkna andraderivatan med (\ref{eqn:fbisscenter}) i punkten \( x_0=7 \)  med samma värden på
\( h  \)  som i Uppgift 1 och jämför med det exakta värdet som kan beräknas på samma sätt som 
i Uppgift 2 fast vi måste lägga till 

\begin{lstlisting}[language=Python]
f_biss = f_prim.diff()
f_biss_x_0 = f_biss.subs(x,7)
print(f_biss_x_0.evalf(16))
\end{lstlisting}

Exakta värdet ska bli 1.806906812550099. Använd även samma värde på \(h\)  som i Uppgift 3
och extrapolera.

\newpage
\section{Uppgift 5}
Nu ska vi ta reda på samtliga nollställen till funktionen:
\begin{equation} \label{eqn:F}
F(x)=-\left(\frac{x}{5}+\cos x\right)
\end{equation}

Börja med att plotta kurvan genom att skriva:
\begin{lstlisting}
  x = ... 
  F = ... 
  pyplot 
\end{lstlisting}



\newpage
\section{Uppgift 6}
Problemet i (\ref{eqn:F}) kan lösas med Newton-Raphsons metod:
\begin{equation} \label{eqn:newton_rhapson}
x_{n+1}=x_n-\frac{F(x_n)}{F'(x_n)}
\end{equation}


\textbf{Svara på följande frågor i rapporten}\\
Algoritmen i (\ref{eqn:newton_rhapson}) konvergerar inte alltid. När blir det problem i (\ref{eqn:newton_rhapson})? 
Vad händer då man startar med \( x_0 = 3  \)? Bestäm samtliga nollställen
till (\ref{eqn:F}) med hjälp av (\ref{eqn:newton_rhapson}) och anteckna antalet
iterationer. Testa med 6 st olika toleranser för lösningen.


\newpage
\section{Uppgift 7}
Problemet i (\ref{eqn:F}) kan även lösas med sekantmetoden:
\begin{equation} \label{eqn:sekantmetoden}
    x_{n+2}=x_{n+1}-\frac{(x_{n+1}-x_n)F(x_{n+1})}{F(x_{n+1})-F(x_n)}
\end{equation}


Bestäm samtliga nollställen till (\ref{eqn:F}) med hjälp av 
(\ref{eqn:sekantmetoden}) och anteckna antalet iterationer. 
Testa med 6 st olika toleranser för lösningen.



\textbf{Svara på följande frågor i rapporten}\\
Fundera vilka fördelar och nackdelar som finns med de olika metoderna 
(\ref{eqn:newton_rhapson}) och (\ref{eqn:sekantmetoden}) och blev det 
någon skillnad i antalet iterationer?


\newpage
\section{Uppgift 8}
I den sista uppgiften ska vi lösa problemet i (\ref{eqn:F}) med Sympys ``nsolve'':

\begin{lstlisting}[language=Python]
x = sp.Symbol('x')
print(sp.nsolve(-x/5-sp.cos(x)), start))
\end{lstlisting}

Bestäm samtliga nollställen.












\end{document}
